\chapter{Testing \& Benchmarking}

\section{Introduction}
Software quality assurance is a fundamental pillar in the lifecycle of any robust digital product, ensuring that the application meets both functional requirements and non-functional performance benchmarks. This is particularly critical for applications operating in the travel sector, where users heavily rely on the platform in real-time while navigating unfamiliar environments, often under varying network conditions. Software validation in mobile applications goes beyond simply checking whether code compiles; it encompasses a holistic evaluation of the user experience, application state management, resilience to adverse conditions, and performance under load. This chapter outlines the testing strategies and software validation methodologies employed throughout the Kemora project, focusing on automated End-to-End (E2E) integration testing on the frontend and rigorous load benchmarking on the backend API. 

\section{Software Validation in Mobile Applications}
Theoretical software validation focuses on the systematic process of evaluating a system to determine whether it satisfies its specified requirements. In mobile applications, validation introduces unique challenges due to device fragmentation, varying screen sizes, differing OS versions, and limited hardware resources. The validation process generally follows a V-Model or Agile continuous testing paradigm, where verification (building the product right) and validation (building the right product) occur simultaneously. Key dimensions of mobile software validation include functional testing to ensure core journeys operate smoothly, usability testing for accessibility and responsiveness, and performance testing to prevent memory leaks and minimize battery consumption. 

\section{Testing Methodology: Unit Testing vs. Integration Testing}
A comprehensive testing strategy requires a layered approach, typically modeled as a "Testing Pyramid."
\subsection{Unit Testing}
Unit testing forms the base of the pyramid, focusing on isolating individual functions, classes, or components and verifying their correctness in a controlled environment. In Kemora, unit tests were employed to validate business logic, data parsing, and state management reducers independently of the UI. Mocking frameworks are heavily utilized at this level to simulate external dependencies such as API clients and local databases, allowing tests to execute rapidly and reliably.

\subsection{Integration and E2E Testing}
While unit testing ensures the internal consistency of individual components, it cannot guarantee that these components interact correctly. Integration testing bridges this gap by evaluating the system as a complete, integrated entity. End-to-End (E2E) testing takes this a step further by simulating real user behavior through the UI, interacting with the application exactly as a human would. This includes tapping buttons, entering text, scrolling, and navigating between screens, validating both the frontend rendering and the integration with the live backend environment. For the Flutter application, the \texttt{integration\_test} package was utilized to simulate a real user driving the application on a physical device or emulator.

\section{End-to-End (E2E) Integration Testing}
The Kemora application utilizes Flutter's native \texttt{integration\_test} package to run E2E tests. This package allows tests to run directly on a target device (simulator, emulator, or physical hardware), capturing the exact performance and rendering behavior of the application.

\subsection{Test State Management and Resetting}
A critical challenge in E2E testing is ensuring test isolation. Tests must not share state, as residual data from a previous test can lead to flaky or inaccurate results. The Kemora test suite programmatically addresses this before initializing the application lifecycle. Specifically, the test begins by explicitly clearing \texttt{SharedPreferences} and resetting the \texttt{TokenStorage} singleton:
\begin{verbatim}
final prefs = await SharedPreferences.getInstance();
await prefs.clear();
TokenStorage.instance.clearTokens();
\end{verbatim}
This explicit state clearing guarantees that the application launches in a clean, unauthenticated state, simulating a fresh installation and ensuring the onboarding and authentication flows can be rigorously validated without interference from previous test runs.

\subsection{Flutter Driver Commands}
The automated user interactions are facilitated through the \texttt{WidgetTester} API, which provides a suite of commands to inspect and manipulate the UI tree:
\begin{itemize}
    \item \textbf{\texttt{tester.pump()} and \texttt{tester.pumpAndSettle()}:} Used extensively to drive the rendering pipeline. \texttt{pumpAndSettle()} repeatedly requests frames until there are no more scheduled animations, which is crucial when waiting for navigation transitions, API responses, or dialog animations to complete.
    \item \textbf{\texttt{tester.tap()}:} Simulates a physical screen tap on a resolved UI widget.
    \item \textbf{\texttt{tester.enterText()}:} Injects string payloads into \texttt{TextField} or \texttt{TextFormField} widgets, utilized heavily in the authentication and AI trip generation forms.
    \item \textbf{\texttt{tester.dragUntilVisible()} and \texttt{tester.ensureVisible()}:} Emulates a user scrolling gesture or forces layout resolution, necessary for interacting with elements that are off-screen within a \texttt{CustomScrollView} or \texttt{ListView}.
\end{itemize}

Finding widgets is achieved using the \texttt{CommonFinders} class, which allows the driver to resolve UI elements via \texttt{find.text()}, \texttt{find.byType()}, \texttt{find.byIcon()}, \texttt{find.widgetWithText()}, \texttt{find.byKey()}, or \texttt{find.byTooltip()}.

\subsection{Test Orchestration}
The E2E suite is orchestrated via a custom PowerShell script (\texttt{run\_e2e\_tests.ps1}). This script automates the environment setup:
\begin{enumerate}
    \item It ensures the local .NET API server is running to serve backend requests.
    \item It starts the ChromeDriver to facilitate communication with the Flutter web/desktop targets.
    \item It executes the \texttt{flutter drive} command, pointing to the \texttt{app\_test.dart} integration test file.
\end{enumerate}

\subsection{User Journey Coverage}
The integration test suite (\texttt{app\_test.dart}) programmatically executes a comprehensive simulation of the core user lifecycle:
\begin{itemize}
    \item \textbf{Onboarding \& Authentication Flow:} The test verifies the presence of the onboarding screen, skips it, and navigates to the registration form. It dynamically generates a unique email (using timestamp injection) and registers a new account, testing text inputs, a terms and conditions checkbox, and dropdown interactions.
    \item \textbf{Exploration Flow:} After successful login, the test validates the Home Screen layout, asserts the presence of a personalized welcome message, and waits for place feeds to load.
    \item \textbf{Social Interaction:} The driver navigates to the Social tab, opens the post creation view via the floating action button, submits a new post ("Hello from E2E integration test!"), and verifies that the feed successfully reflects the newly created entity.
    \item \textbf{AI Trip Generation:} The most complex sequence involves navigating to the Trip tab, inputting "Cairo" as a destination, and triggering the "Generate Magic Itinerary" API. The driver utilizes extended \texttt{pumpAndSettle()} timeouts (up to 15 seconds) to accommodate the LLM generation latency, subsequently asserting that the itinerary renders correctly in the UI. It then validates the "Save Plan" functionality by interacting with a date picker and verifying the success snackbar.
    \item \textbf{Profile \& Logout:} Finally, the test navigates to the Profile tab, validates the previously generated user credentials are displayed, logs out, and asserts the successful return to the Login screen.
\end{itemize}

\section{Backend Load Testing \& Benchmarking}
To guarantee that the .NET 10 API could handle concurrent user traffic—especially during peak travel seasons—a comprehensive load test was conducted.

\subsection{Methodology}
A custom asynchronous Python benchmarking script (\texttt{benchmark.py}) utilizing the \texttt{aiohttp} library was developed. The script targets the core read-heavy data retrieval endpoint \texttt{GET /api/v1/places} (the versioned public places list, which is anonymously accessible and exercises the caching and serialization layers). Concurrency is throttled by a single \texttt{asyncio.Semaphore} instantiated once per run, so the configured number of in-flight requests is genuinely enforced. The load parameters were set to simulate moderate-to-heavy traffic:
\begin{itemize}
    \item \textbf{Total Requests:} 200 consecutive requests.
    \item \textbf{Concurrency Level:} 50 simultaneous connections, throttled via \texttt{asyncio.Semaphore}.
\end{itemize}

\subsection{Benchmark Results}
The benchmark was executed against the local .NET 10 development server (Kestrel, SQL Server backend) on July 1, 2026. All 200 requests completed successfully with a 100\% success rate (HTTP 200), confirming the stability of the stack under moderate concurrency.

\begin{table}[H]
\centering
\begin{tabular}{|l|l|}
\hline
\textbf{Metric} & \textbf{Recorded Value} \\ \hline
Total Requests & 200 (all HTTP 200) \\ \hline
Error Rate & 0.00\% \\ \hline
Average Latency & 103.79 ms \\ \hline
Median (P50) Latency & 108.63 ms \\ \hline
95th Percentile (P95) Latency & 131.63 ms \\ \hline
99th Percentile (P99) Latency & 132.32 ms \\ \hline
\end{tabular}
\caption{API Load Benchmark Results (200 requests, concurrency 50)}
\label{tab:benchmark}
\end{table}

\subsection{Analysis of Results}
The results show a \textbf{median (P50) latency of 108.63 ms} with a tight distribution: the P99 (132.32 ms) sits only $\sim$24 ms above the median, and the average (103.79 ms) closely tracks the median. This narrow spread indicates \emph{predictable, stable} performance rather than erratic tail-latency spikes—no request was an order-of-magnitude outlier.

The 0\% error rate under 50-way concurrency confirms that the Kestrel thread pool and the SQL Server connection pool were never exhausted at this load level; the CLR injected worker threads fast enough to keep the request queue drained. The absolute latency floor (P50 $\approx$ 109 ms) reflects the combined cost of Kestrel routing, controller execution, JSON serialization of the places payload, and any cache lookups. In a production environment behind a robust reverse proxy (e.g., Nginx) or a cloud load balancer with HTTP/2 connection multiplexing, the per-request connection-setup overhead would be amortized, further reducing this floor.

Overall, the benchmark confirms that the backend architecture handles moderate concurrent load robustly and without degradation. Chapter 7 extends this analysis with escalating multi-tiered stress tests to characterize the system's behavior at the limits of the local hardware.
